\SourcePage{754}{757}
repeat the summation over $n$, now using (148.14) for $d\sigma_n$.
Summation by means of (148.21), with $f=d_x$, and the fact that
$\langle d_x\rangle=0$ give
\begin{equation}
 d\sigma_r=\left(\frac{2e}{\hbar v}\right)^2\langle d_x^2\rangle\frac{do}{\vartheta^2}.
 \tag{148.24}
\end{equation}
It is instructive to compare this result with the small-angle elastic-scattering
cross-section (139.5).  The latter is independent of $\vartheta$, whereas the
inelastic cross-section into a solid-angle element $do$ increases as
$1/\vartheta^2$ when $\vartheta$ decreases.

For angles $v_0/v\ll\vartheta\ll1$, and hence $qa_0\gg1$, the second and
third terms within the braces in (148.23) are small.  We then simply obtain
\begin{equation}
 d\sigma_r=Z\left(\frac{2e^2}{mv^2}\right)^2\frac{do}{\vartheta^4},
 \tag{148.25}
\end{equation}
which is Rutherford scattering by the $Z$ atomic electrons, with exchange
neglected.  Recall that the differential elastic-scattering cross-section
(139.6) is proportional to $Z^2$, rather than to $Z$.

Finally, angular integration gives the total cross-section $\sigma_r$ for
inelastic scattering through arbitrary angles and with arbitrary atomic
excitation.  Proceeding exactly as in the calculation of $\sigma_n$ in
(148.18), we find
\begin{equation}
 \sigma_r=8\pi\left(\frac{e}{\hbar v}\right)^2\langle d_x^2\rangle
 \ln\left(\beta\frac{v\hbar}{e^2}\right).
 \tag{148.26}
\end{equation}

\subsection*{Problems\footnote{Atomic units are used throughout these problems.}}

\noindent\textbf{1.} Determine the angular distribution, for
$v^{-2}\ll\vartheta\ll1$, of the inelastic scattering of fast electrons by a
hydrogen atom in its ground state.

\SolutionHeading For hydrogen, the third term in the braces in (148.23) is
absent, while the atomic factor $F(q)$ was calculated in the problem following
\S~139.  Its substitution gives
\[
 d\sigma_r=\frac4{v^4\vartheta^4}\left[1-\left(1+\frac{v^2\vartheta^2}{4}\right)^{-4}\right]do.
\]

\noindent\textbf{2.} Find the differential cross-section for collisions of
electrons with a ground-state hydrogen atom that excite the $n$th discrete
level, where $n$ is the principal quantum number.

\SolutionHeading The matrix elements are conveniently evaluated in parabolic
coordinates.  Choose the $z$ axis along $\mathbf q$; then
$e^{i\mathbf q\mathbf r}=e^{iqz}=e^{iq(\xi-\eta)/2}$.  The ground-state wave
function is $\psi_{000}=\pi^{-1/2}e^{-(\xi+\eta)/2}$.  Matrix elements are
nonzero only for transitions to states with $m=0$.  The wave functions of
these states are
\[
 \psi_{n_1n_20}=\frac1{\sqrt{\pi n^2}}\exp\!\left(-\frac{\xi+\eta}{2n}\right)
 F\!\left(-n_1,1,\frac\xi n\right)F\!\left(-n_2,1,\frac\eta n\right)
\]

\SourcePage{755}{758}
with $n=n_1+n_2+1$.  The required matrix elements are represented by the
integrals
\[
 \langle n_1n_20|e^{i\mathbf q\mathbf r}|000\rangle=
 \int_0^\infty\!\!\int_0^\infty
 \exp\!\left(i\frac q2(\xi-\eta)\right)\psi_{000}\psi_{n_1n_20}
 \frac{\xi+\eta}{4}\,2\pi\,d\xi\,d\eta.
\]
The integrations use the formulas collected in \S~f of the Mathematical
Appendix.  Their evaluation yields
\[
 |\langle n_1n_20|e^{i\mathbf q\mathbf r}|000\rangle|^2
 =2^8n^6q^2\frac{[(n-1)^2+(qn)^2]^{n-3}}{[(n+1)^2+(qn)^2]^{n+3}}
 [(n_1-n_2)^2+(qn)^2].
\]
All states having the same $n_1+n_2=n-1$ have equal energies.  Summing over
all possible values of $n_1-n_2$ at fixed $n$, and inserting the result into
(148.9), gives the required cross-section:
\[
 d\sigma_n=\frac{2^{11}\pi}{v^2}n^7\left[\frac{n^2-1}{3}+(qn)^2\right]
 \frac{[(n-1)^2+(qn)^2]^{n-3}}{[(n+1)^2+(qn)^2]^{n+3}}\frac{dq}{q}.
\]

\noindent\textbf{3.} Determine the total cross-section for excitation of the
first excited state of the hydrogen atom.

\SolutionHeading Integrate
\[
 d\sigma_2=\frac{2^8\pi}{v^2}\frac{dq}{q(q^2+9/4)^5}
\]
over $q$ from $q_{\min}=(E_2-E_1)/v=3/8v$ to $q_{\max}=2v$, retaining only
terms of the highest order in $v$.  Elementary integration gives\footnote{The
cross-section can likewise be calculated for arbitrary $n$.  Numerical
calculation also gives the total inelastic-scattering cross-section for
hydrogen: $\sigma_r=(4\pi/v^2)\ln(v^2/0.160)$.  Of this result, excitation of
discrete-spectrum states and ionization account, respectively, for
$\sigma_{\text{exc}}=(4\pi/v^2)\cdot0.715\ln(v^2/0.45)$ and
$\sigma_{\text{ion}}=(4\pi/v^2)\cdot0.285\ln(v^2/0.012)$.}
\[
 \sigma_2=\frac{2^{18}\pi}{3^{10}v^2}\left(\ln4v-\frac{25}{24}\right)
 =\frac{4\pi}{v^2}\cdot0.555\ln\frac{v^2}{0.50}.
\]

\noindent\textbf{4.} Find the cross-section for ionization of a ground-state
hydrogen atom when the secondary electron is emitted in a specified direction.
The secondary electron's energy is small compared with the primary electron's
energy, so exchange effects are immaterial (H.~Massey, C.~Mohr, 1933).

\SolutionHeading The initial atomic wave function is
$\psi_0=\pi^{-1/2}e^{-r}$.  In the final state the atom is ionized.  Denote the
wave vector of the emitted secondary electron by $\boldsymbol\kappa$, so that
its energy is $\kappa^2/2$.  This state is described by
$\psi_{\boldsymbol\kappa}^{(-)}$ from (136.9), whose ``outgoing'' part at
infinity consists only of a plane wave propagating in the
$\boldsymbol\kappa$ direction.  The function
$\psi_{\boldsymbol\kappa}^{(-)}$ is normalized to a delta function in
$\boldsymbol\kappa/2\pi$-space.  The cross-section calculated from it is
therefore referred to $d^3\kappa/(2\pi)^3$, or to
$\kappa^2d\kappa\,do_{\kappa}/(2\pi)^3$, where $do_{\kappa}$ is the
solid-angle element for the secondary electron's direction.  Thus
\[
 d\sigma=\frac{4k'\kappa^2}{(2\pi)^3kq^4}
 |\langle\boldsymbol\kappa|e^{-i\mathbf q\mathbf r}|0\rangle|^2
 do\,do_{\kappa}\,d\kappa
\]

\SourcePage{756}{759}
where $do$ is the solid-angle element for the scattered electron, and
\[
 \langle\boldsymbol\kappa|e^{-i\mathbf q\mathbf r}|0\rangle
 =\int\psi_{\boldsymbol\kappa}^{(-)*}e^{-i\mathbf q\mathbf r}\psi_0\,dV
 =\frac{e^{-\pi/2\kappa}\Gamma(1-i/\kappa)}{\pi^{1/2}}I,
\]
\[
 I=\left\{-\frac\partial{\partial\lambda}\int
 \exp(-i\mathbf q\mathbf r-i\boldsymbol\kappa\mathbf r-\lambda r)
 F\!\left(\frac i\kappa,1,i(\kappa r+\boldsymbol\kappa\mathbf r)\right)
 \frac{dV}{r}\right\}_{\lambda=1}.
\]
Carry out the integration in parabolic coordinates with the $z$ axis along
$\boldsymbol\kappa$.  Let $\varphi$ be measured from the plane
$(\mathbf q,\boldsymbol\kappa)$:
\[
\begin{split}
 I=\biggl\{&-\frac12\frac\partial{\partial\lambda}
 \int_0^\infty\!\!\int_0^\infty\!\!\int_0^{2\pi}
 \exp\!\biggl[-\frac i2q(\xi-\eta)\cos\gamma
 +iq\sqrt{\xi\eta}\sin\gamma\cos\varphi\\
 &\hspace{28mm}-\frac\lambda2(\xi+\eta)-\frac i2\kappa(\xi-\eta)\biggr]
 F(i/\kappa,1,i\kappa\xi)\,d\varphi\,d\xi\,d\eta\biggr\}_{\lambda=1}.
\end{split}
\]
Here $\gamma$ is the angle between $\boldsymbol\kappa$ and $\mathbf q$.
The integrations over $d\varphi\,d\eta$ are readily performed by substituting
$\sqrt\eta\cos\varphi=u$ and $\sqrt\eta\sin\varphi=v$, after which
\[
 \frac I{2\pi}=\left\{\frac\partial{\partial\lambda}\int_0^\infty
 \exp\!\left[\frac{-q^2\sin^2\gamma+\lambda^2+(\kappa+q\cos\gamma)^2}
 {2[i(\kappa+q\cos\gamma)-\lambda]}\xi\right]
 \frac{F(i/\kappa,1,i\kappa\xi)\,d\xi}{i(\kappa+q\cos\gamma)-\lambda}
 \right\}_{\lambda=1}.
\]
The remaining integral follows from (f.3) with $\gamma=1$ and $n=0$.

The subsequent calculations are lengthy but elementary.  They lead to the
following cross-section:
\[
\begin{split}
 d\sigma={}&\frac{2^8k'\kappa
 [q^2+2q\kappa\cos\gamma+(\kappa^2+1)\cos^2\gamma]}
 {\pi kq^2[q^2+2q\kappa\cos\gamma+1+\kappa^2]^4}\\
 &\times\frac{\exp\!\left\{-\dfrac2\kappa
 \arctan\dfrac{2\kappa}{q^2-\kappa^2+1}\right\}}
 {[(q+\kappa)^2+1][(q-\kappa)^2+1](1-e^{-2\pi/\kappa})}
 do\,do_{\kappa}\,d\kappa.
\end{split}
\]
Integration over all directions of emission of the secondary electron is
elementary and gives, at the specified emitted-electron energy $\kappa^2/2$,
the distribution of scattering directions
\[
 d\sigma=\frac{2^{10}k'\kappa[q^2+(1/3)(1+\kappa^2)]
 \exp\{-\frac2\kappa\arctan[2\kappa/(q^2-\kappa^2+1)]\}}
 {kq^2[(q+\kappa)^2+1]^3[(q-\kappa)^2+1]^3(1-e^{-2\pi/\kappa})}
 do\,d\kappa.
\]
For $q\gg1$ this expression has a sharp maximum near $\kappa=q$.  Close to
the maximum,
\[
 d\sigma=\frac{2^5}{3\pi\kappa^4}\frac{d\kappa\,do}{[1+(q-\kappa)^2]^3}.
\]
Integrating over
$do=2\pi q\,dq/k^2\approx(2\pi\kappa/k^2)d(q-\kappa)$ gives
$8\pi\,d\kappa/k^2\kappa^3$, in agreement, as required, with the first term
of (148.17).
